\chapter*{KATA PENGANTAR}
\addcontentsline{toc}{chapter}{KATA PENGANTAR}

Puji dan syukur penulis panjatkan ke hadirat Tuhan Yang Maha Esa atas segala rahmat dan karunia-Nya, sehingga penulis dapat menyelesaikan Laporan Tugas Akhir yang berjudul "Analisis dan Evaluasi Prediksi Harian PM10 Jakarta Menggunakan Model Hibrida Random Forest-ARIMA Berbasis Rekayasa Fitur". Laporan ini disusun sebagai salah satu syarat untuk menyelesaikan program pendidikan Sarjana pada Program Studi Sistem dan Teknologi Informasi, Sekolah Teknik Elektro dan Informatika, Institut Teknologi Bandung.

Penulis menyadari bahwa penyelesaian Tugas Akhir ini tidak terlepas dari bimbingan, dukungan, serta bantuan dari berbagai pihak. Oleh karena itu, penulis mengucapkan terima kasih dan penghargaan yang sebesar-besarnya kepada:

\begin{enumerate}
    \item Allah SWT atas segala rahmat dan karunia-Nya, serta semua pihak yang telah mendoakan dan mendukung kelancaran Tugas Akhir ini.
    \item Ibu Dr. Fetty Fitriyanti Lubis, S.T., M.T., selaku Dosen Pembimbing atas segala bimbingan dan arahan, serta Bapak I Gusti Bagus Baskara Nugraha, selaku Ketua Program Studi atas dukungan akademik yang diberikan.
    \item Keluarga tercinta, khususnya kedua orang tua dan adik, atas doa, perhatian, serta dukungan moral dan material yang senantiasa diberikan kepada penulis.
    \item Seluruh teman baik dan rekan kerja atas kolaborasi, bantuan, dan kebersamaan yang saling menguatkan dalam menyelesaikan masa studi ini.
\end{enumerate}

\bigskip

Penulis menyadari bahwa Laporan Tugas Akhir ini masih memiliki batasan dan kekurangan. Oleh karena itu, kritik dan saran yang membangun sangat diharapkan untuk perbaikan di masa mendatang. Akhir kata, semoga laporan ini dapat memberikan manfaat bagi pembaca serta menjadi referensi bagi penelitian selanjutnya di bidang pemodelan kualitas udara.

\begin{flushright}
    Bandung, \today\\
    
    Penulis \\
    \vspace{2cm}
    Muhammad Rafi Dhiyaulhaq
\end{flushright}